% -----------------------------------------------------------------------------
% CEC 2022 results in the paper's results-section format.
% This is an input fragment, not a standalone LaTeX document.
% Sources: continuous_benchmark/resultados/run_dtw_job_19070/
%          continuous_benchmark/resultados/run_ddtw_job_19066/
% -----------------------------------------------------------------------------

\subsection{Results in Continuous Optimization: IEEE CEC2022 Suite}
\label{sec:results_cec2022}

The continuous optimization benchmark was evaluated on the IEEE CEC2022
real-parameter test suite, comprising functions F1--F12 at dimension $D=10$.
Each function was evaluated using 31 independent runs and a computational
budget of $T_{\max}=1{,}000$ iterations. The completed campaigns correspond
to the DTW and DDTW variants stored in
\texttt{run\_dtw\_job\_19070} and \texttt{run\_ddtw\_job\_19066},
respectively. Since the benchmark is a minimization problem, lower objective
values indicate better performance.

To report the optimization optima requested for the two framework variants,
Table~\ref{tab:cec2022_results} presents the best objective value observed
among the 31 runs, rather than the arithmetic mean. The column $f^\star$
contains the known CEC reference optimum, while the two gap columns quantify
the relative distance between the best observed value and that reference:
\[
\mathrm{Gap}(\%) =
100\,\frac{f_{\mathrm{best}}-f^\star}{|f^\star|}.
\]
Bold font identifies the lower best value when DTW and DDTW differ; a zero
gap indicates that at least one run reached the reference optimum.

\begin{table}[htbp]
\centering
\caption{Best objective values attained by the DTW and DDTW frameworks on the IEEE CEC2022 benchmark suite ($D=10$, $N=31$ runs).}
\label{tab:cec2022_results}
\resizebox{\textwidth}{!}{
\begin{tabular}{lccccc}
\toprule
\textbf{Function} & \textbf{Optimum ($f^\star$)} & \textbf{DTW best} & \textbf{DDTW best} & \textbf{DTW gap (\%)} & \textbf{DDTW gap (\%)} \\
\midrule
\textbf{F1}  & \textbf{300.00}  & \textbf{300.000} & \textbf{300.000} & 0.000 & 0.000 \\
\textbf{F2}  & \textbf{400.00}  & \textbf{400.011} & \textbf{400.011} & 0.003 & 0.003 \\
\textbf{F3}  & \textbf{600.00}  & \textbf{600.011} & 600.018 & \textbf{0.002} & 0.003 \\
\textbf{F4}  & \textbf{800.00}  & 808.000 & \textbf{806.000} & 1.000 & \textbf{0.750} \\
\textbf{F5}  & \textbf{900.00}  & \textbf{900.000} & \textbf{900.000} & 0.000 & 0.000 \\
\textbf{F6}  & \textbf{1800.00} & 1846.429 & \textbf{1830.400} & 2.579 & \textbf{1.689} \\
\textbf{F7}  & \textbf{2000.00} & \textbf{2000.501} & 2001.495 & \textbf{0.025} & 0.075 \\
\textbf{F8}  & \textbf{2200.00} & \textbf{2201.214} & 2207.625 & \textbf{0.055} & 0.347 \\
\textbf{F9}  & \textbf{2300.00} & \textbf{2529.284} & \textbf{2529.284} & 9.969 & 9.969 \\
\textbf{F10} & \textbf{2400.00} & \textbf{2500.258} & 2500.281 & \textbf{4.177} & 4.178 \\
\textbf{F11} & \textbf{2600.00} & \textbf{2600.000} & \textbf{2600.000} & 0.000 & 0.000 \\
\textbf{F12} & \textbf{2700.00} & \textbf{2860.680} & \textbf{2860.680} & 5.951 & 5.951 \\
\bottomrule
\end{tabular}
}
\end{table}

As reported in Table~\ref{tab:cec2022_results}, both framework variants
reached the reference optimum in at least one run on F1, F5, and F11. DTW
obtained the lower best value on F3, F7, F8, and F10, whereas DDTW obtained
the lower best value on F4 and F6. The variants tied at the displayed
precision on F1, F2, F5, F9, F11, and F12. Thus, the best-run comparison
indicates a landscape-dependent balance between the two alignment variants
rather than universal dominance by one of them.

\begin{itemize}
    \item On the unimodal function \textbf{F1}, both variants reached the
    reference optimum of $300.00$. The same exact optimum was also reached on
    \textbf{F5} and \textbf{F11}.
    \item On the hybrid functions, DDTW produced the closest best run on
    \textbf{F4} and \textbf{F6}, with gaps of $0.750\%$ and $1.689\%$,
    respectively. DTW produced the closest best run on \textbf{F7} and
    \textbf{F8}, with gaps of $0.025\%$ and $0.055\%$.
    \item The composition functions \textbf{F9} and \textbf{F12} remained the
    most difficult cases. Both variants reached $2529.284$ on F9, which is
    $9.969\%$ above its reference optimum, and $2860.680$ on F12, with a
    remaining gap of $5.951\%$.
\end{itemize}

These values describe the best outcome found in 31 stochastic repetitions and
should not be interpreted as typical performance across all runs. The complete
run distributions are therefore analyzed separately using per-function
Friedman tests, Holm-adjusted Wilcoxon comparisons, and the direct paired
DTW--DDTW test in Section~\ref{sec:cec2022-dtw-statistics}. That analysis
found no statistically significant direct difference between the two alignment
variants after Holm correction.

The numerical values were obtained from
\texttt{continuous\_benchmark/resultados/run\_dtw\_job\_19070/resumen\_global.csv}
and
\texttt{continuous\_benchmark/resultados/run\_ddtw\_job\_19066/resumen\_global.csv}.

